% Pollis New Methods: survival analysis (Cox model with Breslow's ties) — draft 1 (11 October 2026)
% Compile with: latexmk -pdf nm_survival.tex && latexmk -c
% Shared look for the Pollis New Methods documents (same fonts and colours as the brochure and the workflow handouts).
\documentclass[10pt,a4paper]{article}

\usepackage[a4paper,left=18mm,right=18mm,top=15mm,bottom=17mm,footskip=8mm]{geometry}
\usepackage[T1]{fontenc}
\usepackage[default,semibold]{sourcesanspro}
\usepackage{textcomp}
\usepackage[italic]{mathastext}
\usepackage{amsmath}
\usepackage{microtype}
\usepackage[table]{xcolor}
\usepackage{tikz}
\usepackage[most]{tcolorbox}
\usepackage{fontawesome5}
\usepackage{tabularx,booktabs}
\usepackage{enumitem}
\usepackage{fancyhdr}
\usepackage[hidelinks]{hyperref}

% ---------------------------------------------------------------- colours (as in the brochure)
\definecolor{navy}{HTML}{1C3A5E}
\definecolor{pblue}{HTML}{3B78B5}
\definecolor{amber}{HTML}{F2B33D}
\definecolor{deepamber}{HTML}{D9902A}
\definecolor{pale}{HTML}{EEF3F9}
\definecolor{ink}{HTML}{2A2F36}
\definecolor{logonavy}{HTML}{173D63}
\definecolor{logodark}{HTML}{C98A1E}
\definecolor{logoamber}{HTML}{F2B446}
\definecolor{logocream}{HTML}{FEEFD2}
\definecolor{soft}{HTML}{5E6773}
\definecolor{agree}{HTML}{2F7D4F}
\definecolor{differ}{HTML}{B23B3B}
\color{ink}

\setlength{\parindent}{0pt}
\setlength{\parskip}{4pt}
\setlist{nosep,leftmargin=*,itemsep=2pt}
\renewcommand{\arraystretch}{1.15}

\pagestyle{fancy}
\fancyhf{}
\renewcommand{\headrulewidth}{0pt}
\fancyfoot[L]{\footnotesize\color{soft}\textbf{Pollis}\enspace New Methods\enspace·\enspace\nmtheme}
\fancyfoot[R]{\footnotesize\color{soft}\thepage}
\newcommand{\nmtheme}{}
\newcommand{\nmdate}{11 October 2026}

% The Pollis logo, redrawn so it prints sharply at any size
\newcommand{\pollislogo}[1]{%
  \begin{tikzpicture}[x=#1/512,y=-#1/512,baseline=(current bounding box.center)]
    \fill[logonavy] (0,0) rectangle (512,512);
    \fill[logodark] (36,422) -- (104,256) -- (172,422) -- cycle;
    \fill[logoamber] (104,422) -- (203,180) -- (302,422) -- cycle;
    \fill[logocream] (203,422) -- (339,89) -- (476,422) -- cycle;
  \end{tikzpicture}}

% Top of the first page: logo, name; then the kicker, title and a one-line promise
\newcommand{\nmhead}[3]{%
  \renewcommand{\nmtheme}{#1}%
  \thispagestyle{fancy}%
  \begin{minipage}[c]{13mm}\pollislogo{11mm}\end{minipage}%
  \begin{minipage}[c]{\dimexpr\linewidth-13mm\relax}
    {\color{navy}\Large\bfseries Pollis}\enspace{\color{soft}New Methods}\hfill
    {\color{deepamber}\small\bfseries\MakeUppercase{#1}}
  \end{minipage}\par\vspace{7mm}
  {\color{navy}\fontsize{22}{26}\selectfont\bfseries #2}\par\vspace{4pt}
  {\color{soft}\large #3}\par\vspace{5mm}}

\newcommand{\nmsection}[1]{\par\vspace{5pt}{\color{navy}\large\bfseries #1}\par\vspace{1pt}}
\newcommand{\nmsub}[1]{\par\vspace{2pt}{\color{deepamber}\bfseries #1}\par}
\newcommand{\gt}{\,{\color{deepamber}\bfseries ›}\,}
\newcommand{\menupath}[1]{{\bfseries\color{navy}#1}}
\newcommand{\fn}[1]{\texttt{\small #1}}
% Citations: \nmcite{key}{Breiman (2001)} links to the entry \nmref{key}{...} in the References
\newcommand{\nmcite}[2]{\hyperlink{ref:#1}{\color{pblue}#2}}
\newcommand{\nmref}[2]{\raisebox{1.1em}[0pt][0pt]{\hypertarget{ref:#1}{}}#2\par}
\newcommand{\nmdoi}[1]{\href{https://doi.org/#1}{\color{pblue}doi:#1}}
% Where a notebook is in Pollis: its title in Browse Methods, and its file
\newcommand{\browse}[2]{\menupath{Compose \gt New Methods \gt Browse Methods}: \emph{#1} (\fn{#2.ipynb})}

% Key facts at the top: source, missing in Julia, notebook, result
\newcommand{\factsheet}[4]{%
  \begin{tcolorbox}[enhanced, colback=pale, colframe=pale, boxrule=0pt, sharp corners,
    left=7pt, right=7pt, top=5pt, bottom=5pt, fontupper=\small]
  \begin{tabularx}{\linewidth}{@{}>{\leavevmode\bfseries\color{navy}}p{27mm}>{\leavevmode}X@{}}
  Tested against & #1\\
  Missing in Julia & #2\\
  Notebook & #3\\
  Result & #4
  \end{tabularx}
  \end{tcolorbox}}

% Boxes
\newtcolorbox{pbox}[2][]{enhanced, colback=pale, colframe=amber, boxrule=0pt,
  leftrule=3pt, sharp corners, left=7pt, right=7pt, top=6pt, bottom=6pt,
  fontupper=\small\raggedright, before upper={{\color{navy}\normalsize\bfseries #2}\par\smallskip}, #1}
\newtcolorbox{nbox}[2][]{enhanced, colback=navy, colframe=navy, boxrule=0pt,
  sharp corners, left=8pt, right=8pt, top=7pt, bottom=7pt, coltext=white,
  fontupper=\small\raggedright, before upper={{\color{amber}\normalsize\bfseries #2}\par\smallskip}, #1}

% Value tables: published beside Julia
\newenvironment{valuetable}[1]{%
  \par\small\tabularx{\linewidth}{#1}\toprule}{\bottomrule\endtabularx\par}
\newcommand{\ok}{{\color{agree}\faCheck}}
\newcommand{\no}{{\color{differ}\faTimes}}

% End of the document: credit and licence
\newcommand{\nmcredit}{%
  \par\vfill
  \begin{minipage}[b]{0.78\linewidth}\footnotesize\color{soft}
  Methods written and tested with Claude Code (Anthropic) in Pollis, from the published papers and their data.\par
  \textcopyright\ 2026 Antonio Saragga Seabra\enspace·\enspace Licensed under \href{https://creativecommons.org/licenses/by-sa/4.0/}{CC BY-SA 4.0}\enspace·\enspace Draft, \nmdate
  \end{minipage}\hfill
  \begin{minipage}[b]{0.18\linewidth}\raggedleft\pollislogo{9mm}\par\vspace{2pt}{\footnotesize contact@pollis.org}\end{minipage}}


\begin{document}

\nmhead{Survival Analysis}
  {The Cox model with Breslow's ties, and its predicted survival}
  {The Mayo Clinic model for primary biliary cirrhosis, reproduced from the 1989 paper.}

\factsheet
  {\nmcite{dickson1989}{Dickson, Grambsch, Fleming, Fisher and Langworthy (1989)}, \emph{Hepatology}: Table 2 (both models) and Table 4 (baseline
   survival, years 1--7). The 418 Mayo patients as \nmcite{fleming1991}{Fleming and Harrington (1991)} published them (PollisDatasets \fn{pbc}).}
  {Breslow's rule for tied deaths, and the predicted survival curve of a patient from a fitted Cox model. Survival.jl uses
   Efron's rule and does not predict survival curves.}
  {\browse{The Cox model with Breslow's ties, and its predicted survival}{cox}}
  {the initial model and all seven years of predicted survival to every printed digit; 41 of 53 values in all. With
   Survival.jl's Efron ties, only 1 to 3 decimals.}

\nmsection{The problem}
When several patients die on the same day, the partial likelihood of \nmcite{cox1972}{Cox (1972)} needs a rule for the tie. Survival.jl uses the rule of
\nmcite{efron1977}{Efron (1977)}. Much of the published literature, and the SAS, BMDP and early S software it was computed with, used that of \nmcite{breslow1974}{Breslow (1974)},
so its coefficients cannot be reproduced in Julia. A model that cannot be reproduced cannot be checked, updated or applied
to new patients with confidence. And a clinical model is used through its predictions: ``what is this patient's chance of
surviving five years?'', which Survival.jl does not compute.

\nmsection{The method}
\nmsub{Breslow's ties}
All patients who die on the same day share one risk set: each death contributes
$\eta_i - \log \sum_{j \in R(t_i)} e^{\eta_j}$, with $\eta = X\beta$. \fn{cox} maximises it by Newton's method with exact
derivatives from ForwardDiff, and returns coefficients, standard errors from the inverse information, Wald chi-squares and
the likelihood ratio chi-square.
\nmsub{Predicted survival}
A patient with risk score $R = X\beta$ has $S(t, R) = S_0(t)^{\exp(R - c)}$, $c$ the mean risk score. The $S_0$ of \nmcite{kalbfleisch1980}{Kalbfleisch and
Prentice (1980)} is a product, over the days with deaths, of factors $\alpha$, each the root of
$\sum_{\text{deaths}} w/(1-\alpha^w) = \sum_{\text{at risk}} w$ with $w = \exp(R - c)$; \fn{survival} finds each root by bisection.

\nmsection{What the tests show}
\begin{valuetable}{@{}lXl@{}}
\textbf{Test} & \textbf{What is compared} & \textbf{Agreement}\\\midrule
Survival.jl & data without ties: Breslow is then Cox & 6 decimals or more\\
Table 2, initial & 312 trial patients: all 20 values & every printed digit\\
Table 2, final & 418 patients & last digit in 6 values (see below)\\
Table 2 & likelihood ratio chi-square & 199.13 vs 199.2 (open)\\
Table 4 & $S_0(t)$, years 1 to 7 & every printed digit\\
Efron (Survival.jl) & Table 2, for contrast & only 1 to 3 decimals\\
\end{valuetable}
The contrast in the last row is the point: on the same data, with the same covariates, the only difference between a fit
that reproduces the paper and one that does not is the rule for ties.

\nmsub{The final model, in the last digit}
Log bilirubin 0.8706 vs 0.8707, log prothrombin time 2.379 vs 2.380, and four chi-squares by 0.01 or 0.02. The paper
estimated missing values for 16 declined patients with regressions fitted to the 90 complete ones (p.\,3); the public file
lacks only two prothrombin times, so the regression here is simpler. Setting those two times to anything from 9 to 13 seconds
moves the coefficients in the third decimal, gaps of this size. The initial model, which needs no estimated values, agrees
throughout.

\nmsub{The paper's two patients}
A 52-year-old with bilirubin 0.5, albumin 4.5, prothrombin time 10.1 and no edema: the paper gives a 96\% chance of
surviving 5 years. One with bilirubin 13.9, albumin 2.8, prothrombin time 13.8 and edema treated with diuretics: 43\% of
surviving 1 year. The notebook draws both curves from the final model.

\nmsection{Still open}
\begin{itemize}
  \item \textbf{The likelihood ratio chi-square}, 199.13 against 199.2, although every coefficient and standard error it comes
    from agrees. Not explained.
  \item \textbf{The one-sample log-rank test} the paper used to validate the model on 106 patients (p.\,5) is not included:
    its groups use a biopsy variable that is not in the public data, and their sizes are not given, so there are no expected
    results to test against.
\end{itemize}

\nmsection{Further considerations}
\begin{itemize}
  \item \textbf{Ties matter more than they seem.} Survival data recorded in days, months or years are full of ties; the
    choice between Breslow and Efron is a convention of the software, and a reader of an older paper has to know which was used.
  \item \textbf{Upstream.} A \fn{ties} option and a \fn{predict} for survival curves in Survival.jl would make every
    Breslow-era model in the literature reproducible in Julia. Tables 2 and 4 are ready-made tests.
  \item \textbf{Next.} Time-varying covariates (the Mayo model was later updated for repeated visits), stratification,
    and confidence bands for $S(t, R)$, each with a published table to test it.
\end{itemize}

\nmsection{References}
{\footnotesize
\nmref{breslow1974}{Breslow, N. (1974). Covariance analysis of censored survival data. \emph{Biometrics} 30(1), 89--99.}
\nmref{cox1972}{Cox, D.\,R. (1972). Regression models and life-tables. \emph{Journal of the Royal Statistical Society, Series B} 34(2), 187--220.}
\nmref{dickson1989}{Dickson, E.\,R., Grambsch, P.\,M., Fleming, T.\,R., Fisher, L.\,D. and Langworthy, A. (1989). Prognosis in primary biliary cirrhosis: model for decision making. \emph{Hepatology} 10(1), 1--7. \nmdoi{10.1002/hep.1840100102}}
\nmref{efron1977}{Efron, B. (1977). The efficiency of Cox's likelihood function for censored data. \emph{Journal of the American Statistical Association} 72(359), 557--565.}
\nmref{fleming1991}{Fleming, T.\,R. and Harrington, D.\,P. (1991). \emph{Counting Processes and Survival Analysis}. Wiley. Appendix D.}
\nmref{kalbfleisch1980}{Kalbfleisch, J.\,D. and Prentice, R.\,L. (1980). \emph{The Statistical Analysis of Failure Time Data}. Wiley.}
}

\nmcredit
\end{document}
